\documentclass[12pt]{article}
\usepackage{amsmath, amssymb, geometry, graphicx, tikz}
\usepackage{actuarialangle}
\usetikzlibrary{patterns}
\usetikzlibrary{calc}
\geometry{margin=1in}
\setlength{\parskip}{1em}
\setlength{\parindent}{0pt}

\begin{document}

\begin{center}
    \Large \textbf{Useful Formulas}
\end{center}

\section{Accumulation Functions \& Interest rate conversions}

\subsection{Accumulation Functions}

(1.) Effective interest rate from $t_1$ to $t_2$: $i_{[t_1,t_2]} = \dfrac{a(t_2)-a(t_1)}{a(t_1)}$

(2.) Effective discount rate from $t_1$ to $t_2$:  $d_{[t_1,t_2]} = \dfrac{a(t_2)-a(t_1)}{a(t_2)}$

(3.) Simple Interest accumulation function: $a(t) = 1+rt$ (where $r = i_{[0,1]}$)

(4.) Simple Discount accumulation function: $a(t) = \frac{1}{1-rt}$ (where $r = d_{[0,1]}$)

(4.) Compound Interest accumulation function: $a(t) = (1+i)^t$ (where $i = i_{[t,t+1]}$ for any $t$)

(5.) Discount function: $v(t)=\dfrac{1}{a(t)}$ and for compound interest $v = v(1)= \dfrac{1}{1+i}$

\subsection{Compound Interest Rate Conversions}

(1.) $1+i = \left(1+\frac{i^{(m)}}{m}\right)^m $ where $i^{(m)}$ is the nominal interest rate compounded $m$ times per year

(2.) $1-d = \left(1-\frac{d^{(m)}}{m}\right)^m $ where $d^{(m)}$ is the nominal discount rate compounded $m$ times per year

(3.) $d = \dfrac{i}{1+i}$ and $i = \dfrac{d}{1-d} = \dfrac{d}{v}$

(4.) Constant force of interest or nominal rate convertible continuously: 
\vspace{-.5cm}
    
    \[\delta = \ln(1+i) \iff e^{\delta}=1+i\]


(5.) Force of interest at time $t$: $\delta_t = \dfrac{a'(t)}{a(t)} = \dfrac{d}{dt}\left(\ln(a(t)\right) \iff a(t) = e^{\displaystyle \int_0^t \delta_r \, dr}$


\section{Annuity Formulas}

(1.) Present value of an annuity immediate (value at time $t=0$ of payments of $\$1$ at times $t=1,2,...,n$) \[a_{\actuarialangle{n}i} = \dfrac{1-v^n}{i}\]

(2.) Value at time $t$ of payments of $\$1$ at times $t = 1,2,...,n$: $a_{\actuarialangle{n}i} (1+i)^t$. When $t =n$: \[s_{\actuarialangle{n}i}=a_{\actuarialangle{n}i }(1+i)^n\]

(3.) Present value of an annuity due: $\ddot{a}_{\actuarialangle{n}i   } = a_{\actuarialangle{n}i  } (1+i)$ but if you insist on knowing the formula we have \[\ddot{a}_{\actuarialangle{n}i} = \dfrac{1-v^n}{d}\]

(4.) Perpetuities: $a_{\actuarialangle{\infty}i} = \dfrac{1}{i}$ and $\ddot{a}_{\actuarialangle{\infty}i} = \dfrac{1}{d}$

(5.) Annuity with geometric payments (growth rate $g$):\[\text{PV}=\dfrac{1-\left(\frac{1+g}{1+i}\right)^n}{i-g}\]

(6.) Annuity with arithmetic payments (starting payment $P$ and additive growth $Q$) \[(I_{P,Q}a)_{\actuarialangle{n}i} = Pa_{\actuarialangle{n}i   }+\frac{Q}{i}(a_{\actuarialangle{n}i}-nv^n)\]

(7.) Perpetuity with increasing payments:
\vspace{-.5cm}

\begin{itemize}
    \item Geometric progression: $\dfrac{1}{i-g}$
    \item Arithmetic progression: $\dfrac{P}{i}+\dfrac{Q}{i^2}$
\end{itemize}

(8.) Annuities with payments less frequent ($k$ times with effective rate $I$) than the given interest period (with rate $i$) \[a_{\actuarialangle{n}I} = \dfrac{a_{\actuarialangle{n}i}}{s_{\actuarialangle{k}i}}\]

(8.) Annuities with payments more frequent than the given interest period (with rate $i$) \[a^{(m)}_{\actuarialangle{n}i} = a_{\actuarialangle{n}i} \cdot \dfrac{i}{i^{(m)}} \hspace{1cm} \text{and} \hspace{1cm} \ddot{a}^{(m)}_{\actuarialangle{n}i} = a_{\actuarialangle{n}i} \cdot \dfrac{i}{d^{(m)}}\]

(9.) Continuously paying annuities at a rate of $f(t)$: $\overline{a}_{\actuarialangle{n}i} = \int_0^n f(t) v(t) dt$

\section{Loans}

(1.) Outstanding loan balance at time $t$ (directly after any payment at time $t$) of a loan $L$ if every payment is \textit{exactly} $P$
\vspace{-.5cm}
\begin{itemize}
    \item Retrospective: $B(t) = L(1+i)^t-Ps_{\actuarialangle{t}i}$ (Loan moved forward $-$ FV of payments made)
    \item Prospective: $B(t) = Pa_{\actuarialangle{n-t}i}$ (PV of future payments)
\end{itemize}

(2.) Recursive outstanding loan balances at time: $B(t+1)= B(t)(1+i)-\text{Pmt}(t)$

(3.) Interest in payment at time $t$: $\text{In}(t) = iB(t-1)$ 

(4.) Principal in payment at time $t$: $\text{Pr}(t) = \text{Pmt}(t)-\text{In}(t)$

(5.) Principal paid from time $t_1$ to time $t_2$: $B(t_1)-B(t_2)$

(6.) Interest paid from time $t_1$ to time $t_2$: $\sum_{t=t_1}^{t_2} \text{Pmt}(t) - (B(t_1)-B(t_2))$

\section{Bond Formulas}

(1.) Coupon amounts: $Fr = Cg = Gj$ 

(2.) The price of a bond: $P = Fra_{\actuarialangle{n} j} + Cv_j^n$

(3.) The premium-discount formula for a bond:
\vspace{-.5cm}
\begin{itemize}
    \item Premium ($P>C$) = $P-C = C(g-j)a_{\actuarialangle{n}j}$
    \item Discount ($C>P$) = $C-P = C(j-g)a_{\actuarialangle{n}j}$ 
\end{itemize}

(4.) Book value at time $t$ using premium-discount formulas: \[B(t) = C(g-j)a_{\actuarialangle{n-t}j}+C\]

(5.) Interest in coupon at time $t$: \[\text{In}(t)=Cg-C(g-j)v_j^{n-t+1}\]

(6.) Amortization of premium or Accumulation of discount at time $t$: \[\text{Pr}(t)=C(g-j)v_j^{n-t+1}\]


\end{document}
